\section{Hardware-Aware Performance Model}
\label{sec:cost-hypothesis}

The model connects completion iterations to executed work and runtime. Its inputs are the observed completion iteration of each problem, independently measured update costs, and calibrated control overhead. This use of workload shape and device cost follows accelerator performance modeling \citep{yu2021habitat,nvidia2022gpuperformance}. Work-count identities supply the workload inputs; device measurements supply the costs. \Cref{sec:execution-controls} tests the resulting time estimates against separate application runs.

\subsection{Completion iterations determine active width}

Let $n_t$ be the iteration of the first configured verification at which problem $t$ passes, and let $n_{\max}=\max_t n_t$. Before update $k$, the active width is
\begin{equation}
A_k=\bigl|\{t:n_t\geq k\}\bigr|.
\label{eq:active-width}
\end{equation}
These iteration counts include the effect of the verification interval. A more frequent schedule can observe completion earlier and produce a different trajectory.

For the same completion decisions, fixed-width execution uses $Wn_{\max}$ problem-update slots. Compaction uses
\begin{equation}
S_{\mathrm{active}}=\sum_{k=1}^{n_{\max}}A_k=\sum_{t=1}^{W}n_t,
\qquad S_{\mathrm{static}}=Wn_{\max}.
\label{eq:update-count-identity}
\end{equation}
The identity follows by counting each problem once at every iteration before its completion. The removed work is
\begin{equation}
R=Wn_{\max}-\sum_t n_t.
\label{eq:removable-updates}
\end{equation}
As illustrated in \cref{fig:method-loop}, $R$ is the area between the fixed-width rectangle and the active-width trajectory. In the full-width logical-mask control, completed states are frozen while their update computations remain. Indexed execution and in-place skipping can also attain $\sum_t n_t$ expensive problem updates while preserving the original state layout.

A batch has removable work when completed and unfinished problems coexist before an update. Thus variation within a batch matters: identical completion iterations yield $R=0$, even if completion iterations differ across batches. The grouping experiment changes this within-batch variation using a fixed set of problems.

\subsection{Kernel savings must cover control cost}

Let $c_k(a,\xi_k)$ denote the measured cost of update $k$ at width $a$. The state $\xi_k$ specifies support shape, data type, kernel configuration, and other conditions of the comparison. The compacted runtime is
\begin{equation}
\begin{split}
T_{\mathrm{active}}={}&\sum_k c_k(A_k,\xi_k)
+H_{\mathrm{screen}}+H_{\mathrm{verification}}\\
&+H_{\mathrm{compact}}+H_{\mathrm{other}}.
\end{split}
\label{eq:time-accounting}
\end{equation}
The $H$ terms account for screening, verification, compaction, and remaining timed control work. Write their sum as $H_{\mathrm{active}}$, and let $\Delta H=H_{\mathrm{active}}-H_{\mathrm{fixed}}$ be the incremental control cost relative to fixed-width execution.

For a comparison with the same completion decisions and matched update states, subtracting the two runtime accounts gives
\begin{equation}
\begin{split}
T_{\mathrm{fixed}}-T_{\mathrm{active}}
={}&\sum_{k=1}^{n_{\max}}
[c_k(W,\xi_k)-c_k(A_k,\xi_k)]\\
&-\Delta H.
\end{split}
\label{eq:time-difference}
\end{equation}
Compaction reduces runtime when
\begin{equation}
\sum_k[c_k(W,\xi_k)-c_k(A_k,\xi_k)]>\Delta H.
\label{eq:active-gate}
\end{equation}
The width sweep measures the update-cost term, and component measurements quantify the additional control work. This equation explains the break-even condition for the measured execution path.

\subsection{Calibration and prediction}

A runtime estimate for a new batch uses independently calibrated costs $\widehat c(a,\xi)$ and overheads $\widehat H$. Given a completion trace, its conditional speedup estimate is
\begin{equation}
\widehat S=\frac{n_{\max}\widehat c(W,\xi)+\widehat H_{\mathrm{fixed}}}
{\sum_k\widehat c(A_k,\xi)+\widehat H_{\mathrm{active}}}.
\label{eq:conditional-prediction}
\end{equation}
Calibration fixes the device, precision, support shape, tile configuration, state layout, and tail path. Costs are specific to the execution strategy: indexed access, entry skipping, and physical compaction can have different update and control costs even for the same active width.

For solver-time validation, we measure update and screen costs at every active width and profile initialization, verification, and release processing in separate calibration runs. The estimate is
\begin{equation}
\begin{split}
\widehat T={}&\widehat H_{\mathrm{init}}+\sum_k\widehat c_{\mathrm{update}}(A_k,\xi)\\
&+\sum_{j\in\mathcal J}\left[\widehat c_{\mathrm{screen}}(A_j,\xi)
+\widehat h_{\mathrm{verify}}(v_j)
+\widehat h_{\mathrm{release}}(d_j)\right],
\end{split}
\label{eq:calibrated-runtime}
\end{equation}
where $\mathcal J$ indexes checks, $v_j$ counts verifier candidates, and $d_j$ counts released problems. Verification uses a calibrated per-candidate cost; release processing uses a calibrated event cost, with separate costs for empty checks. The observed trajectory and event counts determine the estimate, while the target run's total time is reserved for evaluation. We report signed error $100(\widehat T/T-1)$ in \cref{sec:execution-controls}. Target-group transfer holds the evaluated targets out of calibration and retains the shared source input. Pipeline estimates additionally require consumer costs under the same configuration.

\subsection{GPU scheduling makes width cost stepwise}

For the packed Triton update, width $a$ and support size $n$ produce
\begin{equation}
B(a,n)=a\left\lceil\frac{n}{b_m}\right\rceil
\label{eq:block-grid}
\end{equation}
thread blocks, where $b_m$ is the row tile size. GPU scheduling executes blocks in waves whose capacity depends on occupancy \citep{nvidia2022gpuperformance,nvidia2024cuda}. A model of this cost is
\begin{equation}
c(a;n,\xi)\approx g_{n,\xi}\!\left(
\left\lceil\frac{B(a,n)}{Q_\xi}\right\rceil\right),
\label{eq:wave-cost}
\end{equation}
where $Q_\xi$ is the resident block capacity for the device and kernel, and $g_{n,\xi}$ maps wave count to elapsed time. Registers and shared memory affect $Q_\xi$.

Reducing width within a scheduling wave can leave update time nearly unchanged. Larger problems launch more blocks per lane, so removing a lane can reduce the number of waves earlier. \Cref{sec:width-results} measures this relation on A100 and uses it to interpret the application results.
